% Uncompiled proof handoff. Intended as an input to a document using amsmath,
% amssymb, amsthm and theorem/lemma/proposition environments.
\section{The microscopic decorated-AKLT realization, in explicit coordinates}
\label{urgent:aklt-realization}

This module supplies the microscopic input to the two-port completely
positive comparison theorem. It does not assume a positive patch gap.
Instead it constructs the patch matrices and gives an explicit, very
small, positive lower bound for their least nonzero eigenvalues. All
interactions below are norm-one orthogonal projectors. The underlying
core graph is finite and simple, has no isolated vertices in the regions
under discussion, and has maximum degree at most $\Delta\geq2$.

\subsection{Spin one, Kraus matrices, and the Choi constants}

Put
\[
 X=\begin{pmatrix}0&1\\1&0\end{pmatrix},\qquad
 Y=\begin{pmatrix}0&-i\\i&0\end{pmatrix},\qquad
 Z=\begin{pmatrix}1&0\\0&-1\end{pmatrix}.
\]
The symmetric two-qubit physical basis is
$|00\rangle,(|01\rangle+|10\rangle)/\sqrt2,|11\rangle$.
Represent its three coefficient matrices by
$S_+=E_{00}$, $S_0=(E_{01}+E_{10})/\sqrt2$, and $S_-=E_{11}$.
Contracting the second virtual leg with the invariant antisymmetric
form $J=\left(\begin{smallmatrix}0&1\\-1&0\end{smallmatrix}\right)$
and multiplying the local tensor by $\sqrt{2/3}$ gives
\[
 A_+=\sqrt{2/3}\,E_{01},\qquad
 A_0=-Z/\sqrt3,\qquad A_-=-\sqrt{2/3}\,E_{10}.
\]
The physical change of orthonormal basis
\[
 K_1=(A_+-A_-)/\sqrt2=X/\sqrt3,\quad
 K_2=-i(A_++A_-)/\sqrt2=Y/\sqrt3,\quad
 K_3=-A_0=Z/\sqrt3
\]
is unitary on the three Kraus labels. Thus it preserves the CP map.
Changing the normalization of a virtual singlet or reversing its
orientation changes only nonzero scalar normalizations or virtual
unitaries; the displayed convention fixes these choices explicitly.

For every complex matrix $B=\left(\begin{smallmatrix}a&b\\c&d\end{smallmatrix}\right)$,
direct multiplication gives
\[
 XBX=\begin{pmatrix}d&c\\b&a\end{pmatrix},\quad
 YBY=\begin{pmatrix}d&-c\\-b&a\end{pmatrix},\quad
 ZBZ=\begin{pmatrix}a&-b\\-c&d\end{pmatrix}.
\]
Consequently, without a Hermiticity assumption,
\[
 \Phi(B):=\sum_{j=1}^3K_jBK_j^*
 =\frac23\operatorname{tr}(B)I-\frac13B.
\]
Let $\Psi(B)=\operatorname{tr}(B)I/2$ and $\Pi=\mathrm{id}-\Psi$.
Then $\Psi^2=\Psi$, $\Pi^2=\Pi$, $\Psi\Pi=\Pi\Psi=0$, whence
\[
 \Phi^\ell=\Psi+\beta_\ell\Pi,
 \qquad \beta_\ell=(-1/3)^\ell.
\]

Use the unnormalized Choi convention
$J(T)=\sum_{a,b=0}^1E_{ab}\otimes T(E_{ab})$ and
$\Omega=|00\rangle+|11\rangle$. Then
\[
 J(\Psi)=I_4/2,\qquad J(\mathrm{id})=|\Omega\rangle\langle\Omega|,
\]
and the Choi matrix relative to the faithful reference metric is
\[
 J(\Psi)^{-1/2}J(\Phi^\ell)J(\Psi)^{-1/2}
 =(1-\beta_\ell)I_4+4\beta_\ell|\omega\rangle\langle\omega|,
 \qquad \omega=\Omega/\sqrt2.
\]
The one-dimensional space $\mathbb C\omega$ has eigenvalue
$1+3\beta_\ell$; its orthogonal complement has dimension three and
eigenvalue $1-\beta_\ell$. This proves both eigenvalues and multiplicities.
The midpoint of the extreme eigenvalues is $1+\beta_\ell$, irrespective
of parity. Accordingly the centered reference is
\emph{$\widetilde\Psi_\ell=(1+\beta_\ell)\Psi$}, and
\[
 (1-\delta_\ell)\widetilde\Psi_\ell
 \leq_{\rm CP}\Phi^\ell
 \leq_{\rm CP}(1+\delta_\ell)\widetilde\Psi_\ell,
 \qquad
 \delta_\ell=\frac{2|\beta_\ell|}{1+\beta_\ell}
 =\frac2{3^\ell+(-1)^\ell}.
\]
For $\ell\geq2$, the reference is faithful and $0<\delta_\ell<1$.
For $\ell=1$ the lower Choi eigenvalue is zero, explaining the strict
length restriction. The centering scalar must not be silently omitted.

\subsection{The maximal-spin projector as multiplication followed by its adjoint}

For $p\geq0$, write $V_p$ for binary homogeneous polynomials of degree
$p$, and equip its monomial basis
$e_i^{(p)}=x^{p-i}y^i$ with
\[
 \langle e_i^{(p)},e_j^{(p)}\rangle
 =\delta_{ij}\binom pi^{-1}.
\]
This is the Hermitian form inherited from $\operatorname{Sym}^p\mathbb C^2$;
it is invariant under $SU(2)$. For $p,q\geq1$ let
\[
 m_{p,q}:V_p\otimes V_q\longrightarrow V_{p+q},
 \qquad F(x_v,y_v,x_w,y_w)\longmapsto F(x,y,x,y).
\]
Equivalently $m(e_i^{(p)}\otimes e_j^{(q)})=e_{i+j}^{(p+q)}$.
The adjoint is therefore
\[
 m^*e_s^{(p+q)}
 =\sum_{i+j=s}\frac{\binom pi\binom qj}{\binom{p+q}s}
 e_i^{(p)}\otimes e_j^{(q)}.
\]
Vandermonde's identity gives $mm^*=I$. Hence $m^*$ is an isometry and
$P_{p,q}:=m^*m$ is the orthogonal projector onto its image. Its matrix
entries in this nonorthonormal monomial basis are the explicit rationals
\begin{equation}
 (P_{p,q})_{(i,j),(a,b)}
 =\mathbf1_{i+j=a+b}
   \frac{\binom pi\binom qj}{\binom{p+q}{i+j}}.
 \label{urgent:projector-entries}
\end{equation}

To identify the image as the maximal-spin space, only elementary
$SU(2)$ representation facts are needed. The diagonal action on binary
polynomials differentiates to
\[
 H=x\partial_x-y\partial_y,\qquad E=x\partial_y,\qquad F=y\partial_x.
\]
Any nonzero invariant subspace of $V_r$ is invariant under these
operators. The distinct $H$ eigenvalues isolate a monomial by a polynomial
in $H$; repeated $E$ and $F$ then give every monomial. Thus $V_r$ is
irreducible, of spin $r/2$. The multiplication map is equivariant, so its
adjoint image is a copy of $V_{p+q}$. The weight $p+q$ space in
$V_p\otimes V_q$ is one-dimensional. Therefore that copy is unique and
is exactly the maximal total-spin summand. Formula
\eqref{urgent:projector-entries} is the physical norm-one interaction.

\subsection{One edge: the kernel is exactly the bracket ideal}

Let $b_{vw}=x_vy_w-y_vx_w$. Multiplication by $b_{vw}$ maps
$V_{p-1}\otimes V_{q-1}$ injectively into $V_p\otimes V_q$ because
the polynomial ring is an integral domain. Its image lies in $\ker m$
because setting the two pairs of variables equal annihilates the bracket.
Its dimension is $pq$. On the other hand $m$ is surjective, so
\[
 \dim\ker m=(p+1)(q+1)-(p+q+1)=pq.
\]
Equality of dimensions proves
\[
 \ker P_{p,q}=\ker m=b_{vw}(V_{p-1}\otimes V_{q-1}).
\]
This argument uses neither an unproved Clebsch--Gordan decomposition nor
an assumed tensor-network intersection property.

\subsection{Distinct edge brackets are nonassociate primes}

Work in the finite-variable polynomial ring
$R=\mathbb C[x_v,y_v:v\in W]$. For distinct vertices $v,w$, view
$b_{vw}$ as a polynomial in $y_w$ over
$A=\mathbb C[x_v,y_v,x_w,\text{the other variables}]$.
Its coefficients $x_v$ and $-y_vx_w$ are coprime: $x_v$ is a prime
variable and divides neither $y_v$ nor $x_w$. Thus this linear polynomial
is primitive. It is irreducible over the fraction field of $A$, since
its degree is one. Gauss's lemma proves irreducibility in $A[y_w]=R$.
As $R$ is a UFD, it is prime.

Units of $R$ are nonzero complex scalars. The vertex multidegree of
$b_{vw}$ is one at $v,w$ and zero elsewhere. Two brackets can be
associates only if their unordered endpoint sets coincide. In a simple
graph, brackets for distinct edges are therefore nonassociate. Choose
one orientation per edge; reversing it changes only a sign. Simplicity
is essential: parallel-edge brackets would be associates, and a loop
would give the zero polynomial.

It follows by induction using primality that, for a finite internal
edge set $F$,
\[
 \bigl(\forall e\in F,\ b_e\mid f\bigr)
 \quad\Longleftrightarrow\quad
 \prod_{e\in F}b_e\mid f.
\]
For the induction step, a new prime divides $f$ and is nonassociate to
every previous prime, hence does not divide their product and must
divide the remaining quotient.

\subsection{Regional ground space and the actual tensor image}

Let a microscopic region have vertex set $W$, internal edge set $F$,
and full physical degrees $D_v$. Full degrees include edges omitted at
the regional boundary. The physical polynomial space has multidegree
$(D_v)_{v\in W}$. Put $r_v=D_v-\deg_F(v)\geq0$.
Since the local terms are positive orthogonal projectors,
\[
 \ker\sum_{e\in F}P_e=\bigcap_{e\in F}\ker P_e.
\]
The one-edge identity and prime-product argument give the exact space
\begin{equation}
 \ker H_W
 =\left(\prod_{e\in F}b_e\right)
      \bigotimes_{v\in W}V_{r_v}.
 \label{urgent:regional-kernel}
\end{equation}
Indeed a quotient after division by the homogeneous product has exactly
the residual multidegree $(r_v)$; uniqueness of multihomogeneous
components proves this assertion without allowing extra degrees.

The AKLT virtual construction produces exactly this space. An internal
virtual singlet is the antisymmetric tensor with polynomial $b_e$.
At a vertex the physical symmetrization is multiplication of its
$D_v$ linear factors. Therefore contracting all internal singlets and
leaving the boundary virtual tensor free gives, up to fixed nonzero
normalization scalars, multiplication by $\prod_e b_e$. Each boundary
symmetrization $(\mathbb C^2)^{\otimes r_v}\to V_{r_v}$ is surjective:
every monomial is the image of a product of coordinate vectors. Taking
tensor products proves surjectivity onto the residual factor in
\eqref{urgent:regional-kernel}. Correlated boundary tensors are permitted,
so their linear image is the entire displayed tensor product, not merely
its decomposable vectors. This proves the regional tensor-image identity.

For a region comprising whole half-wire blocks, every exposed boundary
wire endpoint has $r_v=1$, and every other vertex has $r_v=0$. Thus the
boundary symmetrization has no kernel in this case. Multiplication by
the nonzero bracket product is injective. This also directly verifies
the expected boundary-space dimension and excludes an additional
boundary-corner kernel.

\subsection{Explicit finite patch types and a positive numerical formula}

A two-block patch joining original core degrees $d,e\in\{1,\ldots,\Delta\}$
contains $\ell$ spin-one vertices on every incident half-wire. Its
physical dimensions, site count, term count and ground-space dimension
are respectively
\[
 D_{d,e}=(d+1)(e+1)3^{\ell(d+e)},\qquad
 N_{d,e}=2+\ell(d+e),\qquad M_{d,e}=1+\ell(d+e),\qquad
 \dim\ker H_{d,e}=2^{d+e-2}.
\]
The last formula follows from the $d+e-2$ exposed endpoint legs in
\eqref{urgent:regional-kernel}. The patch is a tree: the two core stars
are joined by their one central interface. Its type depends only on
$(d,e,\ell)$, up to relabeling. There are at most $\Delta^2$ ordered
types, whether or not every type occurs in a particular graph.

Each $H_{d,e}$ is the sum of the explicitly specified projector matrices
\eqref{urgent:projector-entries}, tensored with identities on other
sites. These matrices are self-adjoint and positive in the physical
weighted inner product. In a Euclidean coordinate basis they are
similar, by the positive diagonal square root of the Gram matrix, to
Hermitian positive matrices. Their rational monomial coordinates
therefore have the same real nonnegative eigenvalues. There is at least
one nonzero interaction, so $H_{d,e}\ne0$ and its rank is positive.

For a fully explicit lower bound define the positive integer
\[
 C_\Delta=\operatorname{lcm}
 \left\{\binom rs:0\leq s\leq r\leq\Delta+2\right\},\quad
 R=(\Delta+1)^2 3^{2\ell\Delta},\quad M=2\ell\Delta+1.
\]
Every bond joins a core of degree $d\leq\Delta$ to degree two, or two
degree-two wire sites. Thus its total polynomial degree is at most
$\Delta+2$ (including $4\leq\Delta+2$). Formula
\eqref{urgent:projector-entries} shows that every entry of
$H_{d,e}$ lies in $C_\Delta^{-1}\mathbb Z$. There is no product of
denominators over bonds: summing matrices with this common denominator
preserves it.

Let the positive eigenvalues of a patch be
$0<\lambda_1\leq\cdots\leq\lambda_r$, where $1\leq r\leq D_{d,e}\leq R$.
The coefficient of the lowest nonzero power in its characteristic
polynomial, in absolute value, is $\prod_{j=1}^r\lambda_j$.
Equivalently this is the $r$th elementary symmetric function of all
eigenvalues, or the sum of the principal $r$-minors. Since the entries
have denominator dividing $C_\Delta$, this coefficient belongs to
$C_\Delta^{-r}\mathbb Z$. It is strictly positive, so
\[
 \prod_{j=1}^r\lambda_j\geq C_\Delta^{-r}.
\]
Each interaction is a norm-one projector, giving
$\lambda_j\leq\|H_{d,e}\|\leq M_{d,e}\leq M$. Hence
\[
 \lambda_1\geq C_\Delta^{-r}M^{1-r}
 \geq \underbrace{C_\Delta^{-R}M^{1-R}}_{=:g_{\mathrm{exp}}(\Delta,\ell)>0}.
\]
This is an explicit finite arithmetic expression. No spectral gap is
assumed in obtaining it. If
\[
 g_*(\Delta,\ell)=\min_{1\leq d,e\leq\Delta}
                  \lambda_{\min}^{+}(H_{d,e}),
\]
then $0<g_{\mathrm{exp}}\leq g_*$; these quantities are not asserted
equal. Exact spectral or interval calculations can substantially improve
this deliberately coarse lower bound, but are unnecessary for existence.

\subsection{Coverage and microscopic comparison}

For a connected region of at least two whole blocks, every internal
microscopic bond lies in a two-block patch. A bond inside block $v$
occurs in precisely the patches associated with included coarse edges
incident to $v$, hence between one and $\Delta$ patches. A bond across
a block interface occurs in its unique interface patch. Therefore
\[
 \Delta H_{\rm micro}\geq\sum_eH_{P_e}
 \geq g_*\sum_e(I-G_e)
 \geq g_{\rm exp}H_{\rm block}.
\]
The regional kernel identity and coverage show that both Hamiltonians
have the same kernel. Consequently a coarse bound
$H_{\rm block}\geq c(I-G)$ gives the actual normalized physical bound
\[
 H_{\rm micro}\geq\frac{c\,g_{\rm exp}}\Delta(I-G).
\]
In particular the independent CP-angle theorem's choice
$\ell=\lceil\log_9(144\Delta)\rceil$, giving $c\geq1/2$, implies
\[
 \operatorname{gap}H_{\rm AKLT}
 \geq \frac{C_\Delta^{-R}(2\ell\Delta+1)^{1-R}}{2\Delta}>0.
\]
The constant is uniform over the stated region family for fixed
$\Delta,\ell$. It is not a degree-uniform estimate and says nothing
about eliminating decoration. This microscopic implication is conditional
only on the separately supplied coarse CP-angle theorem, not on an
unproved microscopic seed gap.
